% Generated by ChimCanhCut Research Platform
% Template: IEEEtran — IEEE Transactions on Neural Networks and Learning Systems / TPAMI / TKDE
% Topic: Federated Learning under Concept Drift: A Unified Perspective on Drift-Aware Learning in Distributed Non-Stationary Environments

\documentclass[journal]{IEEEtran}

\usepackage[utf8]{inputenc}
\usepackage{cite}
\usepackage{amsmath,amssymb}
\usepackage{graphicx}
\usepackage{booktabs}
\usepackage{hyperref}
\usepackage{url}
\usepackage{multirow}
\usepackage{array}

\title{Federated Learning under Concept Drift: A Unified Perspective on Drift-Aware Learning in Distributed Non-Stationary Environments}
\author{[Author Names]}

\begin{document}
\maketitle

\begin{abstract}
% TODO: Enter your abstract via the Paper Metadata panel in the UI.
% Required: 150-250 words covering motivation, methodology, key findings, contributions.
\end{abstract}

\begin{IEEEkeywords}Federated Learning, Concept Drift, Non-IID Data, Distribution Shift\end{IEEEkeywords}

\section{Introduction}
\label{sec:introduction}

Federated learning (FL) is a leading approach for privacy-preserving collaborative machine learning \cite{Collins2025}. The framework allows distributed clients, including Internet of Things (IoT) devices and edge nodes, to jointly train global models without exposing raw local data [CITE:34, CITE:104]. Despite these advantages, real-world distributed environments remain inherently dynamic. Although extensive literature addresses static non-independent and identically distributed (non-IID) data challenges [CITE:90, CITE:91, CITE:92], concept drift—where the underlying data distribution shifts over time—poses an increasingly critical problem \cite{Martin2025}. 

Within distributed networks, concept drift occurs both temporally and spatially across heterogeneous clients, creating what is known as distributed concept drift [CITE:14, CITE:22]. Deploying traditional FL aggregation methods like FedAvg in these nonstationary environments frequently causes severe model degradation and catastrophic forgetting [CITE:48, CITE:78]. Edge and IoT deployments exacerbate this issue, as strict limits on communication bandwidth, computational capacity, and battery life make continuous retraining or complex model maintenance impractical [CITE:73, CITE:101, CITE:122].

In response, researchers have proposed various drift-aware federated learning frameworks. To provide a structured overview of this evolving field, our survey introduces a comprehensive taxonomy categorizing current solutions along three primary dimensions: adaptation strategies, network topologies, and the locus of adaptation. 

First, *adaptation strategies* broadly separate into active and passive methods. Active techniques depend on explicit drift detection mechanisms—including statistical triggers, loss-based monitoring, or dual-schedulers—to initiate reactive model updates [CITE:4, CITE:27, CITE:34, CITE:45]. Passive strategies, by contrast, utilize continuous adaptation techniques like dynamic learning rate adjustments or regularization. These maintain model plasticity during drift phases without relying on computationally expensive statistical detectors [CITE:17, CITE:31, CITE:40, CITE:44]. 

Second, *network topologies* determine how the communication infrastructure manages drift. Solutions range from traditional centralized FL, where a single server orchestrates adaptation, to decentralized (peer-to-peer) and hierarchical architectures. The latter localize drift responses to minimize communication latency and prevent single points of failure [CITE:55, CITE:88]. 

Third, the *locus of adaptation* identifies the primary site of drift mitigation. Client-side adaptation emphasizes localized learning and personalized models [CITE:11, CITE:14], whereas server-side adaptation involves intelligent aggregation and global model clustering [CITE:22, CITE:48]. Hybrid approaches distribute the adaptive workload across both clients and the server, effectively balancing personalization with global generalization [CITE:60, CITE:105].

Within this taxonomic framework, structural model maintenance under drift varies significantly. While certain centralized frameworks adapt a single global model \cite{Casado2021}, other studies argue that distributed concept drift requires multimodel or clustered approaches to effectively capture divergent client distributions [CITE:11, CITE:14, CITE:22]. To navigate the inherent trade-off between accommodating new concepts and retaining past knowledge, researchers increasingly integrate federated continual learning (FCL) techniques. Methods such as experience replay and elastic weight consolidation treat concept drift as a continuous learning problem [CITE:112, CITE:117, CITE:128]. Nevertheless, deploying FCL in edge environments demands careful optimization; the memory and processing overhead of retaining historical knowledge must be directly balanced against hardware limitations.

\begin{figure}[!t]
\centering
\includegraphics[width=\columnwidth]{figures/year_distribution}
\caption{Publication year distribution of included papers.}
\label{fig:year_dist}
\end{figure}

\section{Background and Related Work}
\label{sec:background}

Federated Learning (FL) has become a foundational paradigm for privacy-preserving collaborative intelligence. It allows distributed clients to train machine learning models without exposing raw data \cite{Collins2025}. Despite its success in safeguarding privacy, FL faces well-documented difficulties with spatial data heterogeneity, widely known as the non-independent and identically distributed (non-IID) data problem [CITE:90, 91, 92]. To mitigate the impact of skewed client distributions, researchers have extensively explored techniques like clustered FL, personalized FL, and data augmentation [CITE:94, 99, 102, 105, 107]. Nevertheless, traditional FL frameworks typically assume a static environment; they presume that local data distributions remain stationary over time, even if they differ across clients.

Real-world applications—especially those processing streaming data—rarely operate in static environments. Instead, the statistical properties of input data or target variables frequently shift over time, causing what is known as concept drift [CITE:75, 77, 79]. Centralized machine learning traditionally manages this drift via continuous monitoring and adaptation, relying on distribution-based, performance-based, or unsupervised detectors [CITE:71, 78, 81]. Furthermore, systematic reviews of streaming data classification emphasize a range of active detection algorithms alongside passive adaptation strategies, including dynamic ensemble selection and incremental learning [CITE:55, 63, 74, 83].

When concept drift strikes a federated network, the system must grapple with the compounded challenges of spatial heterogeneity and temporal nonstationarity. This combination creates "distributed concept drift" [CITE:14, 22]. Empirical evidence shows that failing to explicitly address drift in FL can severely degrade global model performance [CITE:33, 42]. Managing these shifts in a distributed setting is further complicated by inherent information asymmetry. Because the central server lacks direct access to local data, it cannot easily determine which specific clients are experiencing drift during the aggregation phase.

To navigate these obstacles, researchers have proposed diverse methodologies. This survey categorizes these approaches into a unified taxonomy structured around three core dimensions: *adaptation strategies*, *network topologies*, and the *locus of adaptation*. First, *adaptation strategies* separate active approaches—which depend on explicit drift detection mechanisms like statistical triggers to launch reactive updates [CITE:34, 43, 45]—from passive ones. Passive strategies employ continuous dynamic adjustments, such as learning rate modulation, to preserve model plasticity without explicit detectors [CITE:17, 31, 40, 44]. Second, *network topologies* dictate the structural response to drift. These range from maintaining a single evolving global model to deploying decentralized or clustered multimodel architectures that capture divergent distributions as new concepts emerge [CITE:1, 11, 14, 21, 22]. Finally, the *locus of adaptation* identifies where primary mitigation takes place: either on the server side (using specialized aggregation weights and momentum) or the client side (through local model regularization or personalized fine-tuning) [CITE:27, 38, 52].

A critical aspect of client-side adaptation involves mitigating catastrophic forgetting, which is the loss of previously acquired knowledge when learning new concepts from nonstationary data. Consequently, Continual Learning (CL) techniques are increasingly integrated into FL frameworks to treat concept drift as a continuous learning problem [CITE:48, 112, 117]. Federated Continual Learning (FCL) frameworks employ strategies like prioritized experience replay, elastic weight consolidation, and knowledge distillation to retain prior concepts while adapting to new ones [CITE:118, 122, 128]. However, integrating these sophisticated FCL mechanisms introdu

\section{Problem Formulation}
\label{sec:problem_formulation}

The problem of Federated Learning (FL) under concept drift emerges at the intersection of spatial data heterogeneity and temporal nonstationarity. In standard FL formulations, a central server—or a network of nodes—coordinates $N$ distributed clients to minimize a global empirical risk function without directly accessing local datasets. Traditional FL implicitly assumes stationary underlying data distributions. However, real-world distributed environments like the Internet of Things (IoT) and edge networks frequently violate this assumption [CITE:34, 39, 73]. 

Concept drift introduces temporal heterogeneity. Specifically, the joint probability distribution of input features and target labels for a given client $i$ at time $t$, denoted as $P_i^t(X, Y)$, changes such that $P_i^t(X, Y) \neq P_i^{t+1}(X, Y)$ [CITE:33, 42, 79]. When this temporal evolution occurs across multiple nodes holding non-independent and identically distributed (non-IID) data [CITE:90, 91, 92], it manifests as *distributed concept drift* [CITE:14, 22]. Solving distributed concept drift requires addressing several interconnected challenges. To provide a unified perspective, we categorize the problem space into three primary dimensions: (1) *drift adaptation strategies*, (2) *the primary locus of adaptation*, and (3) *network topologies*.

The first dimension, **drift adaptation strategies**, addresses how the system identifies and responds to nonstationarity. A fundamental dichotomy exists between active and passive approaches. Active paradigms explicitly monitor statistical triggers, loss metrics, or data distributions to identify drift and initiate reactive model updates [CITE:34, 43, 75, 78]. Frameworks utilizing explicit detection (e.g., FedConD, FLARE) can precisely trigger retraining or model transfers [CITE:34, 39, 45]. Nevertheless, they often assume resource-constrained edge devices can continuously evaluate complex statistical triggers without prohibitive latency overheads—a computationally demanding requirement [CITE:73, 88]. Conversely, passive strategies bypass explicit detection, relying instead on mechanisms like dynamic learning rates or regularization to continuously maintain model plasticity [CITE:17, 31, 40, 44]. Furthermore, evolving local streams leave models susceptible to catastrophic forgetting. Consequently, adaptation frequently intersects with Federated Continual Learning (FCL). This intersection treats drift mitigation as a continuous process that balances plasticity (learning new concepts) and stability (retaining prior knowledge) under strict resource constraints [CITE:48, 112, 117, 122].

The second dimension, the **locus of adaptation**, delineates whether drift mitigation occurs primarily at the client or server level. Client-side methods focus on localized mechanisms—such as personalized learning rates or local memory buffers—to handle drift prior to aggregation [CITE:17, 112]. Server-side techniques involve adjusting aggregation weights, managing model clusters, or routing specialized models to specific clients [CITE:11, 20]. This distinction directly impacts the overall architectural formulation. Maintaining a single global model proves insufficient when distributed drift forces clients into distinct conceptual spaces. To resolve this divergence, multimodel and clustered FL approaches (e.g., FedDrift, FedCCFA) dynamically group clients experiencing similar drifts [CITE:14, 22]. These formulations, however, introduce scalability challenges, including multiple model management, complex merging heuristics, and increased communication overhead. Additionally, privacy constraints render the server largely blind to local data distributions, complicating the accurate identification of drifting clients during server-side aggregation [CITE:4, 27].

The third dimension involves **network topologies**, which

\section{Methodology}
\label{sec:methodology}

\#\# Methodology

We conducted this systematic survey according to the Preferred Reporting Items for Systematic Reviews and Meta-Analyses (PRISMA) 2020 guidelines. This approach provides a transparent and reproducible synthesis of the literature on federated learning (FL) in nonstationary environments. The methodology includes a specific search strategy, strict eligibility criteria, a multistage screening process, and structured data extraction protocols.

\#\#\# 1. Search Strategy
To identify relevant studies at the intersection of FL and concept drift, we searched major academic databases and preprint repositories. The primary sources were **Semantic Scholar** and **arXiv**, with supplementary indexing planned for **IEEE Xplore (pending)**. This combination ensures broad coverage across computer science literature and engineering applications.

We restricted the search period to the years **2015 to 2026**. The 2015 baseline captures early decentralized machine learning paradigms that immediately preceded the formal introduction of FL. Extending the timeframe to 2026 allows for the inclusion of recent advancements and accepted forthcoming publications.

Because terminology varies across subfields, we developed **45 distinct search queries**. Using Boolean operators, these queries intersect core concepts from three primary domains: 
*   *Distributed Learning:* ("federated learning" OR "collaborative learning" OR "decentralized learning")
*   *Nonstationary Environments:* ("concept drift" OR "data shift" OR "non-stationary" OR "temporal shift")
*   *Adaptive Mechanisms:* ("drift-aware" OR "continuous learning" OR "dynamic adaptation")

\#\#\# 2. Inclusion and Exclusion Criteria
Retrieved records were evaluated against predefined inclusion and exclusion criteria to maintain the survey's topical focus and scientific rigor. 

**Inclusion Criteria:**
1.  **Published in peer-reviewed venues or as arXiv preprints:** The rapid pace of AI innovation means preprints frequently offer highly relevant methodologies prior to formal peer review.
2.  **Directly or adjacently related to the survey topic:** Studies must address explicit concept drift in FL or related nonstationary challenges (e.g., continual learning, domain adaptation) that inform drift-aware FL design.
3.  **Published after 2015:** This aligns with the temporal scope of modern FL architectures.
4.  **Abstract available:** This is necessary for the initial semantic screening phase.

**Exclusion Criteria:**
1.  **Duplicate entries:** Removed to prevent redundant analysis and skewed quantitative assessments.
2.  **Non-English papers:** Excluded to ensure the accurate interpretation of complex algorithmic concepts.
3.  **Papers without an accessible abstract:** Excluded because they prevent efficient preliminary relevance evaluation.
4.  **Papers published before 2015:** Filtered out to avoid legacy distributed computing models that do not align with modern privacy-preserving FL constraints.

\#\#\# 3. Study Selection and PRISMA Flow
The study selection process followed a structured, multistage PRISMA pipeline:

*   **Identification:** Executing the 45 search queries across the selected databases yielded **129 records**. 
*   **Screening:** All **129 records** underwent a two-tiered screening process based on titles and abstracts. First, rule-based filtering enforced baseline criteria (e.g., date ranges, language, and duplication), excluding **19 records**. Subsequently, a Large Language Model (LLM)-assisted screening protocol evaluated the semantic alignment of the remaining abstracts with the core topic. This LLM-assiste

\begin{figure}[!t]
\centering
\includegraphics[width=0.85\columnwidth]{figures/prisma_flow}
\caption{PRISMA 2020 flow diagram of the systematic literature search.}
\label{fig:prisma}
\end{figure}

\section{Taxonomy}
\label{sec:taxonomy}

We propose a multidimensional taxonomy to systematically categorize the diverse methodologies addressing nonstationary data in distributed environments. This framework provides a unified perspective by classifying the existing literature along three intersecting axes: Drift Adaptation Strategy, Primary Locus of Adaptation, and Network Topology. Examining how algorithms detect or absorb drift (Strategy), where the computational burden resides (Locus), and how communication infrastructure facilitates model updates (Topology) enables a comprehensive evaluation of design trade-offs in current drift-aware federated systems.

**1. Drift Adaptation Strategy**
The approach to managing concept drift serves as a primary categorization point in the literature. These strategies broadly fall into three groups: Active Detection \& Reactive, Passive Adaptation, and Clustering \& Multimodel.

*Active Detection \& Reactive:* Frameworks in this category use explicit mechanisms to identify concept drift and subsequently trigger model updates. Retraining is typically initiated through statistical tests, dual-schedulers, or historical performance metrics. For instance, the FLARE framework conditionally transfers training data and deploys models based on training behavior and inference statistics \cite{Chow2023}, \cite{Chow2023a}. FLAME similarly applies dynamic thresholding for drift detection \cite{Mavromatis2024}, whereas FL-MalDrift integrates on-device statistical algorithms like ADWIN and DDM \cite{Patel2025}. Although effective at isolating exact drift points, active detection often forces clients to continuously evaluate models on new samples or maintain statistical tracking. Such requirements can introduce significant computational and latency challenges for resource-constrained edge and Internet of Things (IoT) devices.

*Passive Adaptation:* In contrast to explicit detection, passive strategies favor continuous, lightweight adjustments that maintain model plasticity during drift phases. This approach mitigates the computational overhead associated with active triggers. A prevalent technique involves scaling learning rates dynamically based on individual training loss, as demonstrated in Adaptive-FedAVG \cite{Canonaco2021a}, \cite{Canonaco2021} and various client-side adaptation algorithms \cite{Saile2024a}, \cite{Saile2024}. Furthermore, Continual Learning (CL) techniques are frequently integrated into passive FL frameworks to prevent catastrophic forgetting. Rather than treating concept drift as a discrete anomaly detection task, these methods view it as an ongoing learning process, employing regularization or replay buffers (e.g., CDA-FedAvg \cite{Casado2021a}, \cite{Casado2021} and Federated Continual Learning frameworks \cite{Hamedi2025}, \cite{Benameur2025}).

*Clustering \& Multimodel:* This approach addresses variations in data distributions across the network. While standard FL maintains a single evolving global model, multimodel frameworks posit that distributed concept drift often necessitates distinct models to capture divergent client distributions. Algorithms such as FedDrift \cite{Jothimurugesan2022}, \cite{Jothimurugesan2022a} cluster clients based on empirical risk, assigning specialized models to handle different emerging concepts. Other notable examples include FedCCFA, which relies on classifier clustering and feature alignment \cite{Chen2024}, \cite{Chen2024a}, and FedPLC, which implements dynamic cluster adaptation to group clients experiencing similar drift patterns \cite{Zhou2026}, \cite{Zhou2026a}.

**2. Primary Locus of Adaptation**
The responsibility for managing concept drift can be localized to different components within the federated architecture. This choice of locus directly impacts communication overhead, privacy guarantees, and the system's responsiveness to local versus global drift.

*   *Client-Side Adaptation:* In this paradigm, the mechanisms for drift handling are executed locally on edge devices. This approach ensures immediate responsiveness to local data shifts and strict privacy preservation, as raw statistical triggers or replay buffers never leave the device.

\begin{table}[!t]
\renewcommand{\arraystretch}{1.3}
\caption{Proposed Multi-Dimensional Taxonomy}
\label{tab:taxonomy}
\centering
\begin{tabular}{p{3cm}p{9cm}}
\toprule
\textbf{Dimension} & \textbf{Categories} \\
\midrule
\textbf{Drift Adaptation Strategy} & Passive Adaptation, Active Detection \& Reactive, Clustering \& Multi-Model \\
\textbf{Network Topology} & Standard Client-Server FL, Hierarchical / Edge-Assisted FL, Decentralized / Asynchronous FL, Non-Federated Data Streams \\
\textbf{Primary Locus of Adaptation} & Client-Side, Server-Side, Hybrid / Joint, Not Applicable \\
\bottomrule
\end{tabular}
\end{table}

\begin{figure}[!t]
\centering
\includegraphics[width=\columnwidth]{figures/taxonomy_heatmap}
\caption{Heatmap of paper distribution across taxonomy dimensions.}
\label{fig:taxonomy_heatmap}
\end{figure}

\section{Literature Review}
\label{sec:literature_review}

Research on concept drift in federated learning (FL) has largely shifted from tackling static, non-independent and identically distributed (non-IID) data [CITE:91, CITE:92] to handling dynamic, nonstationary environments. As highlighted by recent systematic reviews [CITE:2, CITE:33, CITE:79], traditional FL algorithms suffer significant performance degradation when underlying data distributions shift over time. To synthesize these developments, this review organizes the current literature along three structural dimensions: adaptation strategies (active versus passive), network topologies (centralized, hierarchical, and decentralized), and the primary locus of adaptation (server-side versus client-side).

Regarding adaptation strategies, the literature generally splits into active detection and passive adaptation. Active methods rely on explicit drift detection—such as statistical triggers or performance monitoring—to prompt reactive model updates. Frameworks like FL-MalDrift, for instance, embed statistical detectors (e.g., ADWIN, DDM) directly on devices to monitor local data streams, triggering retraining when distribution shifts occur [CITE:88, CITE:102]. Other examples include FLARE [CITE:34, CITE:43], which uses dual-schedulers driven by training behavior, alongside FedConD [CITE:39, CITE:45] and FLAME \cite{Mavromatis2024}, which rely on adaptive drift detection and dynamic thresholding. A major limitation of these active mechanisms, however, is their assumption that edge devices possess the capacity to continuously evaluate models on incoming samples or sustain complex statistical triggers. Consequently, they often incur substantial computational and latency overheads \cite{ren2025}.

Passive approaches, on the other hand, bypass explicit detection entirely in favor of continuous adaptation. Algorithms like Adaptive-FedAVG [CITE:40, CITE:44] adjust learning rates dynamically based on individual training loss, thereby enhancing model plasticity during drift phases. Because nonstationary data frequently leads to knowledge erosion—commonly known as catastrophic forgetting—researchers increasingly integrate continual learning (CL) techniques into these passive frameworks. This resulting paradigm, federated continual learning (FCL), frames concept drift as an ongoing learning challenge [CITE:112, CITE:117]. For example, CDA-FedAvg [CITE:36, CITE:48] augments standard averaging with continual adaptation, while other frameworks use Prioritized Experience Replay or Elastic Weight Consolidation to preserve knowledge across evolving tasks [CITE:118, CITE:128]. By eliminating the need for explicit detection triggers, passive and FCL mechanisms provide lightweight alternatives highly suited to the resource constraints of IoT environments.

Drift-aware FL architectures also differ substantially based on their underlying network topology. In standard centralized (star) configurations, maintaining a single global model is frequently inadequate if clients encounter divergent concept drifts. As a result, multimodel and clustered federated learning techniques—such as FedDrift [CITE:14, CITE:22], FedCCFA [CITE:11, CITE:21], and FedPLC [CITE:20, CITE:37]—dynamically group clients and allocate specialized models to emerging concepts. Although highly effective for personalized adaptation, centralized clustering introduces scalability bottlenecks, particularly regarding model proliferation and complex merging heuristics. To overcome these limitations, recent studies explore decentralized (peer-to-peer) and hierarchical topologies. Decentralized FL facilitates localized drift adaptation via direct client-to-client communication. This reduces dependence on a central server and permits device neighborhoods to organically adjust to regional concept drifts [CITE:142, CITE:155]. Alternatively, hierarchical models deploy edge servers as intermediaries to aggregate local drifts prior to global updates, striking a balance between rapid local adaptation and global consistency [CITE:160, CITE:162].

Finally, the third dimension concerns the locus of adaptation, which determines whether the primary responsibility for managing drift resides at the server or the client

\begin{figure}[!t]
\centering
\includegraphics[width=\columnwidth]{figures/top_cited}
\caption{Top cited papers in the survey corpus.}
\label{fig:top_cited}
\end{figure}

\section{Critical Analysis}
\label{sec:critical_analysis}

Evaluating the drift-aware federated learning literature exposes distinct methodological divergences and structural challenges. To provide a unified perspective, we assess the state-of-the-art across three core dimensions: adaptation strategies, the locus of adaptation, and network topologies. We then examine current benchmarking practices and emerging fairness considerations.

Adaptation strategies typically fall into two categories: active detection and passive adaptation. Active frameworks, such as FLARE \cite{Chow2023} and FedConD \cite{Chen2021}, rely on explicit drift detection. They frequently integrate on-device statistical detectors like ADWIN or DDM to track drift scores [CITE:8, 43, 45]. Despite their precision, these mechanisms often assume that clients can continuously evaluate models on incoming data streams. This requirement can introduce severe computational and latency overheads in resource-constrained edge and IoT environments [CITE:73, 88]. Passive strategies, such as Adaptive-FedAVG [CITE:40, 44], avoid explicit detection altogether. Instead, they dynamically adjust learning rates based on training loss to improve model plasticity [CITE:17, 31]. While more resource-efficient, these methods must carefully balance plasticity and stability to avoid catastrophic forgetting. To manage this trade-off, Federated Continual Learning (FCL) frameworks increasingly adopt techniques like Elastic Weight Consolidation (EWC) and Prioritized Experience Replay (PER) [CITE:48, 112, 117]. Nevertheless, translating these techniques to distributed settings often creates new communication bottlenecks that demand careful optimization [CITE:118, 128].

The locus of adaptation reveals a clear split between client-side and server-side interventions. Client-side methods focus on localized model updates to address concept shifts independently. Conversely, server-side adaptation generally involves modifying aggregation protocols or managing the overall model architecture. A major debate at the server level centers on whether to maintain a single global model or deploy multimodel and clustered architectures. Frameworks like FedDrift [CITE:14, 22] and FedCCFA [CITE:11, 21] group clients and assign specialized models to distinct emerging concepts [CITE:1, 20]. Although effective for divergent distributions, these clustered approaches often face scalability limits over time, which can lead to high storage demands and complex model-merging procedures. Moreover, privacy-preserving protocols inherently constrain server-side adaptation because the central server lacks direct visibility into local data distributions. As a result, distinguishing between drifting and stable clients during aggregation remains a fundamental challenge, sometimes causing suboptimal updates for clients experiencing no drift [CITE:4, 27].

Regarding network topologies, the vast majority of the literature examines centralized (star-topology) federated networks, where a central server coordinates drift adaptation. In contrast, very few studies address concept drift within decentralized or peer-to-peer federated learning environments [CITE:55, 62]. Without a central aggregator in these decentralized setups, local concept drifts propagate asynchronously across the network. Designing consensus mechanisms capable of detecting and adapting to nonstationary data streams without a central coordinator remains a complex and underexplored research area.

Methodologically, comparative analysis suffers from a lack of unified, standardized benchmarks for nonstationary federated learning [CITE:94, 122]. Many studies rely on proprietary datasets or simplistic synthetic datasets.

\section{Future Directions and Open Problems}
\label{sec:future_directions}

While the proposed taxonomy categorizes the literature by adaptation strategies, network topologies, and the primary locus of adaptation, several underexplored areas remain. Rather than reiterating mechanistic differences among existing methods, this section synthesizes the core findings of the survey to outline actionable research trajectories for advancing distributed nonstationary learning.

**1. Unified Benchmarking Frameworks for Complex Topologies**
The absence of standardized benchmarks for evaluating federated learning under concept drift remains a persistent challenge \cite{Mahdi2025}. Current studies often rely on fragmented, proprietary datasets or simplistic synthetic adaptations, such as artificially corrupted MNIST variants [CITE:94, 122], which complicates rigorous cross-paper comparisons. Addressing this gap requires comprehensive, open-source benchmarking frameworks. To accurately simulate real-world distributed nonstationary environments, these benchmarks must incorporate temporal drift, spatial non-IID data skew, and support for diverse network topologies—including centralized, hierarchical, and decentralized architectures [CITE:92, 54, 131].

**2. Topology-Aware Drift Adaptation Strategies**
Most drift-aware federated learning frameworks assume a standard centralized (star) topology, leaving the dynamics of concept drift in alternative network structures largely unexplored. In decentralized (peer-to-peer) or hierarchical setups, the absence of a global aggregator fundamentally alters drift detection and mitigation [CITE:112, 115]. Consequently, investigating topology-aware adaptation strategies is essential, particularly regarding how concept drift propagates through gossip protocols or edge-server hierarchies. Critical frontiers for scalable nonstationary learning include understanding the impact of network topology on drift synchronization and developing localized consensus mechanisms to manage asynchronous concept shifts \cite{Benameur2025}.

**3. Optimizing the Locus of Adaptation: Client-Side Efficiency vs. Server-Side Attribution**
An ongoing trade-off exists regarding the primary locus of adaptation. Client-side active drift detection mechanisms typically assume devices can continuously evaluate models on streaming samples. However, the resulting computational and memory overheads are often impractical for resource-constrained IoT devices [CITE:34, 88]. Passive adaptation strategies, such as dynamic learning rates, provide lightweight alternatives [CITE:17, 44] but frequently sacrifice detection precision. Conversely, server-side adaptation suffers from limited visibility, as the central server must attribute drift to specific clients without accessing local data distributions. Although techniques like local loss-based detection [CITE:4, 27] or federated fuzzy clustering \cite{Stallmann2024} address this issue, they risk inadvertently exposing client-specific behavioral patterns. Practical deployment will therefore require hybrid frameworks that balance lightweight client-side heuristics with privacy-preserving, server-side drift attribution [CITE:95, 134].

**4. Scalable Knowledge Management in Multimodel Architectures**
To manage distributed concept drift—where clients experience divergent data distribution shifts—recent studies have proposed multimodel and clustered federated learning approaches [CITE:14, 22]. Although these methods effectively assign specialized models to emerging concepts, they introduce significant scalability constraints over extended operational periods. Key challenges include the unbounded proliferation of concept-specific models, suboptimal model merging, and increased storage and communication overheads [CITE:99, 105]. Addressing these issues requires scalable multimodel maintenance strategies, with promising directions including the use of compact model weights for knowledge exchange [CITE:12

\section{Conclusion}
\label{sec:conclusion}

This survey provides a unified perspective on managing concept drift in Federated Learning (FL), moving beyond isolated algorithmic fixes to establish a comprehensive taxonomy. Synthesizing the current state of drift-aware FL, we categorized the literature across three core dimensions: drift adaptation strategies, network topologies, and the primary locus of adaptation. Our analysis highlights a clear shift from localized anomaly detection toward holistic, continuous learning frameworks [CITE:112, 117]. Instead of treating drift as an edge-case disruption, modern approaches increasingly embed plasticity directly into the federated optimization process. This integration ensures that models remain relevant across evolving data streams [CITE:122, 128].

Effective drift mitigation fundamentally requires aligning the adaptation strategy with both the network topology and the processing locus. Centralized FL architectures, for example, typically rely on server-side aggregation adjustments to handle global drift [CITE:4, 27]. Conversely, decentralized and hierarchical topologies require client-side or edge-centric mechanisms to manage localized, asynchronous distribution shifts [CITE:14, 22]. The chosen adaptation locus also dictates how the computational burden is distributed. Client-side methods must balance adaptation efficacy against device constraints [CITE:73, 88]. Meanwhile, server-side strategies face the complex task of distinguishing between drifting clients and inherent non-independent and identically distributed (non-IID) data heterogeneity without compromising privacy [CITE:39, 45].

Although the field has made significant progress, several research avenues demand further exploration. A critical priority for future work is developing unified, standardized benchmarks tailored to nonstationary federated environments \cite{Mahdi2025}. Currently, the reliance on disparate datasets and synthetic drift injections severely limits objective comparative analyses. Validating emerging algorithms will require open-source frameworks capable of simulating realistic, spatio-temporal drift across diverse network topologies. Stronger theoretical foundations are also needed, particularly to establish convergence guarantees and generalization bounds for drift-aware aggregation methods operating under strict privacy constraints [CITE:11, 21].

Scalability remains an ongoing challenge, especially as researchers explore multimodel and clustered architectures to address divergent client distributions [CITE:15, 23]. Maintaining dynamic model repositories or executing continuous client clustering introduces substantial communication and storage overheads [CITE:25, 47]. To reduce the resource footprint of continuous adaptation mechanisms, subsequent studies should investigate more efficient locus-specific techniques. Promising directions include lightweight parameter-efficient fine-tuning (PEFT) at the client level and sparse communication protocols at the server level [CITE:31, 40, 99].

The intersection of model fairness and distributed concept drift also requires deeper investigation. Because local data distributions evolve asynchronously, ensuring equitable model performance across diverse demographic groups or client clusters becomes increasingly complex \cite{Badar2023}. A critical frontier in this domain is the development of fairness-aware adaptation strategies that proactively monitor and mitigate group-specific drift, all without requiring centralized access to sensitive attributes [CITE:15, 44].

Ultimately, enabling FL to robustly navigate nonstationary environments is essential for its practical deployment in dynamic systems. The research community can fully realize the potential of drift-aware machine learning in distributed settings by aligning adaptation strategies with appropriate topologies and processing loci, while simultaneously prioritizing standardized benchmarking, scalable architectures, and fairness.

\bibliographystyle{IEEEtran}
\bibliography{references}

\end{document}